\section{Degenerating saddle transitions: action coordinates}
\label{sec:negative-action}

We use an action coordinate to describe a large negative
source-energy interval. Let \(\beta\) be
defined by \(Y+\log Y=X^2/2+\beta\) for the unperturbed negative source.
The passage consists of two parabolic tails joined by a thin region near
the equator. Its limiting transition is determined by an action integral,
but its curvature vanishes on the reversible coefficient set. Absolute
convergence alone does not control the curvature near that set. We estimate
the nonlinear part of the actual transition relative to the distance to
that set. Keeping the section maps and their inverses in physical
coordinates allows this estimate to be used in
\eqref{eq:intro-displacement}.

\subsection{The equation and the section coordinates}

The local scales below are related to the raw coefficients in
Section~\ref{sec:coefficient-regimes}.

Fix a scale \(t>0\) on each coefficient fibre, keeping it fixed as
the orbit varies. Start with \eqref{ann:source-field} at \(a=-1\),
write \(p=t\Lambda_1\), \(s=t^3\Lambda_2\),
\(\gamma=t^2k\), and put \(\xi=tX\), \(\eta=t^2Y\).
After this phase scaling and a positive change of time, the same
form \eqref{ann:source-field} holds in $(\xi,\eta)$ with the new
coefficients
\begin{equation}\label{neg:source-dictionary}
 (a_{\rm new},p_{\rm new},\gamma_{\rm new},s_{\rm new})
 =(-t^2,t^2\Lambda_1,t^2k,t^2\Lambda_2).
\end{equation}
Thus \(p_{\rm new}=tp\), \(s_{\rm new}=s/t\), and the equation is
\begin{equation}\label{neg:field}
\begin{split}
 \dot\xi&=-\eta-t^2P_\Lambda(\xi)
                +At^2L(\xi)\eta-Ct^4L(\xi)^2\eta,\\
 \dot\eta&=-\eta\{\xi+t^2\xi(2k+\Lambda_2\xi)
                         +Ct^4\Lambda_2L(\xi)\eta\},\\
 P_\Lambda(\xi)&=1-\Lambda_1\xi+3k\xi^2+\Lambda_2\xi^3,
 \qquad L(\xi)=k+\Lambda_2\xi .
\end{split}
\end{equation}
The passive coefficients \(A,C\) range over a fixed compact
neighborhood of \((1,1)\). Physical tails have \(\Lambda_2>0\);
at \(\Lambda_2=0\) we will use limits of normalized maps.

For the unperturbed negative source,
\(Y+\log Y=X^2/2+\beta\). Hence its rescaled energy is
\begin{equation}\label{neg:affine-energy}
 \widehat b=t^2\beta+t^2\log(2t^2),\qquad
 \eta+t^2\log(2\eta)=\xi^2/2+\widehat b.
\end{equation}
Energy-label differences are exactly \(t^2\) times source-label
differences; the action integral below keeps these transverse labels.

We specify the intervals on which the estimate will be used. Choose
compact intervals
\[
 J\Subset\operatorname{int}\widehat J_+\Subset(-\infty,0),
 \qquad \widehat J_-\Subset(-\infty,0),
\]
and a compact coefficient set. For \(b\in\widehat J_+\), put
\(x=\sqrt{-2b}\) and define \(r>0\) by
\begin{equation}\label{neg:action-equation}
 A_-(r;\Lambda)=A_+(x;\Lambda),\qquad
 A_+(x;\Lambda)=\int_0^x\frac{u}{P_\Lambda(u)}\,du,
 \quad A_-(r;\Lambda)=\int_0^r\frac{u}{P_\Lambda(-u)}\,du.
\end{equation}
We require that the entry and exit lie in fixed compact subintervals
of \((0,\infty)\), that \(-r^2/2\) lies in
\(\operatorname{int}\widehat J_-\) with a fixed margin, and that
\begin{equation}\label{neg:rootfree-margin}
 P_\Lambda\ge p_0>0
\end{equation}
on a common neighborhood of every traversed segment \([-r,x]\).
The common neighborhood extends to both sides of the limiting exit.
Monotonicity of \(A_-\) then gives comparison exit cuts with opposite
action differences. Enlarged intervals provide the smaller common
images needed for composition and inversion.

Choose a fixed \(L\) beyond the entry and exit intervals. The positive
height pieces between the two sections \(\xi=\pm L\) and the lips
are required to lie in fixed transverse neighborhoods. Each outward
tail lies in a parabolic corridor and reaches the prescribed raw
section, followed by its actual regular transitions. The divided
finite closing map is included on its side. Its multiplier \(m\)
is held fixed as an independent passive parameter in
\(0\le m\le m_0\) during parameter division, as in
Section~\ref{sec:source-normalization}; the cycle equation uses its
actual value \(m>0\). These requirements concern the whole enlarged
input and output intervals.

There is a useful energy coordinate for all small signed \(k\).
Put
\begin{equation}\label{neg:ideal-energy}
\begin{split}
 \ell&=1-At^2k+Ct^4k^2,\qquad
 \nu_k=\frac{6t^2k}{1+2t^2k},\\
 \mathcal B_{t,k}(\xi,\eta)
 &=-\frac12\left[
 \eta^{-\nu_k}\left(\xi^2+\frac1{3k}
             -\frac{2\ell\eta}{1-4t^2k}\right)-\frac1{3k}\right].
\end{split}
\end{equation}
The apparent pole is removable, with
\begin{equation}\label{neg:ideal-energy-zero}
 \mathcal B_{t,0}=\eta+t^2\log\eta-\xi^2/2 .
\end{equation}
Direct differentiation of \eqref{neg:field} gives
\(\dot{\mathcal B}_{t,k}=0\) when
\(\Lambda_1=\Lambda_2=0\). Shrink the parameters so that
\(1+2t^2k\), \(1-4t^2k\), and \(\ell\) are positive.
At the two fixed cuts use the separate coordinates
\begin{equation}\label{neg:cut-labels}
 b_\pm=\mathcal B_{t,k}(\pm L,\eta)+t^2\log2 .
\end{equation}
Here \(\eta\) stays in a compact positive interval and
\(\partial_\eta\mathcal B_{t,k}\) is bounded above and away
from zero. These coordinates and their inverses are smooth in
\(t^2,k,b_\pm\). They agree with \eqref{neg:affine-energy} at
\(k=0\); away from that set, their relation is part of the actual
section maps.

Let \(S_{t,\lambda,k}\), \(\lambda=(\Lambda_1,\Lambda_2)\),
be the passage from \(b_+\) to \(b_-\), and let
\(\Gamma_\pm\) be the maps from these separate coordinates to
the actual input and divided output heights. For references
\(b_\pm^*\) in the enlarged intervals, set
\begin{equation}\label{neg:tail-normalizations}
\begin{split}
 c_\pm&=\Gamma_\pm'(b_\pm^*)>0,
 \qquad a_\pm=\Gamma_\pm(b_\pm^*)-c_\pm b_\pm^*,\\
 \phi_\pm(b)&=b_\pm^*
          +\frac{\Gamma_\pm(b)-\Gamma_\pm(b_\pm^*)}{c_\pm},
 \qquad \Gamma_\pm=a_\pm+c_\pm\phi_\pm .
\end{split}
\end{equation}
On the respective raw images
\(\mathcal I_\pm=\Gamma_\pm(\operatorname{int}\widehat J_\pm)\)
and normalized images
\(\mathcal U_\pm=\phi_\pm(\operatorname{int}\widehat J_\pm)\), put
\begin{equation}\label{neg:exact-composition}
 \Pi_{t,\Lambda}=\phi_-\circ S_{t,\lambda,k}\circ\phi_+^{-1}
       :\mathcal U_+\longrightarrow\mathcal U_- .
\end{equation}
If necessary the domain is restricted to the smaller input interval
whose image has the prescribed output margin. The estimates below
ensure that \(J\subset\mathcal U_+\) for small parameters.

To see the role of both affine factors, write the actual cycle
equation as \(h-a_{\rm f}-mh_{\rm out}=0\), where
\(a_{\rm f}=a_\theta=C_\theta(0)\) and
\(h_{\rm out}=\mathcal C_{\theta,m}(T_\theta(h))\) is the divided
output height from Lemma~\ref{lem:closing}. Substitution of
\eqref{neg:tail-normalizations} gives, without a change of zeros,
\begin{equation}\label{neg:physical-displacement}
 D(u)=a_+-a_{\rm f}-ma_-+c_+u-mc_-\Pi_{t,\Lambda}(u)=0 .
\end{equation}
Thus \(u\) is affine in the actual incoming height. The two
positive first derivatives may tend to zero at different rates;
neither their equality nor a bounded ratio is required.

The zero count therefore requires curvature control of the completed map
\eqref{neg:exact-composition}. Lemma~\ref{neg:lem-completed-curvature}
expresses this curvature with both section contributions; the relative
comparison below estimates them separately.

\subsection{A relative estimate near the reversible coefficients}

For a fixed reference \(b_*\in J\), write
\(\mathfrak T f=f-f(b_*)-(b-b_*)f'(b_*)\).
On a fixed compact interval the quotient norm satisfies
\[
 \|[f]\|_{C^3/\operatorname{Aff}}
      \asymp\|\mathfrak T f\|_{C^3}\asymp\|f''\|_{C^1} .
\]
This quotient isolates the derivatives needed for the curvature count.
The affine coefficients remain in the actual displacement
\eqref{neg:physical-displacement}.

\begin{theorem}[Relative convergence on the actual sections]
\label{neg:thm-relative}
Fix the enlarged intervals, the entry and exit margins, the tail and
regular-transition neighborhoods described above. There exist
\(t_0,\lambda_0,k_0,C>0\), depending only on these data and the
compact passive coefficient set, such that the following holds.
For
\[
 0<t\le t_0,\qquad
 0<\rho_{\rm id}:=(\Lambda_1^2+\Lambda_2^2)^{1/2}\le\lambda_0,
 \qquad \Lambda_2\ge0,\qquad |k|\le k_0,
\]
put \(\theta_i=\Lambda_i/\rho_{\rm id}\). On the fixed smaller
interval \(J\), the map \eqref{neg:exact-composition} satisfies
\begin{equation}\label{neg:relative-estimate}
 \left\|\frac{[\Pi_{t,\Lambda}-\mathrm{id}]}{\rho_{\rm id}}
           -\theta_1[U_0]-\theta_2[V_0]\right\|_{C^3/\operatorname{Aff}}
 \le C\{t^2(1+|\log t|)+\rho_{\rm id}+|k|\},
\end{equation}
where
\begin{equation}\label{neg:limiting-rows}
 U_0(b)=-\frac23(-2b)^{3/2},\qquad
 V_0(b)=\frac25(-2b)^{5/2} .
\end{equation}
The coefficient \(k\) is independently signed. The constants are
uniform in the two separate reference values, held fixed during
parameter integration, and in \(0\le m\le m_0\). No positive
lower bound on \(\theta_2\) is needed. At \(\Lambda_2=0\) the
assertion concerns the one-sided normalized limits.
\end{theorem}

\begin{proof}
On each outward tail use \eqref{eq:parabolic-tail} with dictionary
\eqref{neg:source-dictionary} and the separate initial label
\eqref{neg:cut-labels}. With \(L\) fixed before shrink, the enlarged
windows have \(\eta\asymp r^2\) and denominator \(\asymp r^2\)
throughout division. Thus the tail proof of
Proposition~\ref{prop:source-normalization} applies with its exact
\(j\asymp r^{\lambda_\gamma}\), kernel
\eqref{eq:normalized-difference}, and bounds
\eqref{eq:N-bound}--\eqref{eq:N-parameter-bound} through one extra
phase order. It gives bounded \(C^3\)-valued intrinsic \(s\)-rows
and \(p\)-rows at \(s=0\), including the moving-end error
\(O_{C^3}(s^{1-\nu})\). These differentiations concern only the tails.

At \(p=s=0\), \eqref{neg:ideal-energy} is a first integral;
implicit differentiation gives identity normalized intrinsic limits.
As in \eqref{ann:exact-source-division}, integrate the intrinsic tails
first in \(p\) at \(s=0\), then in \(s\) at fixed \(p\), holding
the two reference values separately fixed. Since
\(\partial_{\Lambda_1}=t^2\partial_p\)
and \(\partial_{\Lambda_2}=t^2\partial_s\), the intrinsic maps and
inverses differ from the identity by
\(Ct^2(|\Lambda_1|+\Lambda_2)\) in \(C^3\).

At the tested coefficients, each raw gain for \(\Lambda_2>0\) satisfies
\[
 d\le C(t^2\Lambda_2)^{2-\nu}
       \le Ct^2\rho_{\rm id},\qquad 0<\nu<1.
\]
The ordinary observers, including divided closing, therefore add
\(O_{C^3}(d)\) by Lemma~\ref{phys:lem-c3-adapters}. Thus
\begin{equation}\label{neg:single-tail-estimate}
 \|\phi_\pm-\mathrm{id}\|_{C^3}
 +\|\phi_\pm^{-1}-\mathrm{id}\|_{C^3}
       \le Ct^2(|\Lambda_1|+\Lambda_2).
\end{equation}
At \(\Lambda_2=0\) use the one-sided normalized limit. Inverses
are on fixed smaller common images; no exterior coefficient derivative
is needed.

For the core, work in a small independent \((\lambda,k)\)-block with
strict core margin and construct \(S_{t,\lambda,k}\) and its first
\(\lambda\)-derivatives in \(C_b^3\), keeping \(t>0,k\) fixed.
Choose an integer \(M>20\), and cut each lip at
\(y=\eta/t^2=t^M\). On a lip use
\(H=\eta-\xi^2/2\) and \(\ell_\eta=\log\eta\).
Equation \eqref{neg:field} gives exactly
\begin{equation}\label{neg:lip-equation}
 \frac{dH}{d\ell_\eta}=t^2\Phi_\pm,
 \qquad
 \Phi_\pm=
 \frac{\xi\eta(2k+\Lambda_2\xi+AL-Ct^2L^2)
              -\xi P_\Lambda+Ct^2\Lambda_2L\eta^2}
      {\xi\{1+t^2(2k+\Lambda_2\xi)\}
                              +Ct^4\Lambda_2L\eta} .
\end{equation}
The signed root \(\xi^2=2(\eta-H)\) is bounded away from zero
on each lip, as is the denominator. Thus the required derivatives
of \(\Phi_\pm\) in \(H,\lambda,k\) are bounded. The logarithmic
clock has length \(O(1+|\log t|)\), giving a change
\(O(t^2(1+|\log t|))\) through total order five; the positive-height
pieces to \(\pm L\) add \(O(t^2)\). Equation~\eqref{neg:ideal-energy}
therefore gives
\begin{equation}\label{neg:lip-label-error}
 b=H+O_{C^5}\bigl(t^2(1+|\log t|)\bigr),
\end{equation}
At \(k=0\), the error is \(t^2\log(2\eta)\), largest at
\(\eta=t^{M+2}\); the removable continuation in \(k\) has the same
mixed bound through order five. The positive entry is
\(\sqrt{-2b}+O_{C^5}(t^2(1+|\log t|))\); the negative exit has
the corresponding estimate about \(-\sqrt{-2b}\).

In the thin core put \(Q=t^2\log y\). Its exact graph equation is
\begin{equation}\label{neg:thin-core}
 Q_\xi=
 \frac{\xi(1+2t^2k+t^2\Lambda_2\xi)
               +Ct^6\Lambda_2(k+\Lambda_2\xi)y}
      {P_\Lambda(\xi)+\ell_\Lambda(\xi)y}
 =\frac{\xi}{P_\Lambda(\xi)}+R,
 \quad
 \ell_\Lambda=1-At^2L+Ct^4L^2 .
\end{equation}
Until the first return to \(y=t^M\), one has \(0<y\le t^M\).
The root-free margin and direct differentiation give
\begin{equation}\label{neg:core-mixed-estimate}
 |\partial_Q^j\partial_\lambda^\alpha R|
       \le C\{t^2+t^{M-2j}\},\qquad j+|\alpha|\le5,
\end{equation}
with the same bounds for the required \(\xi\)-derivatives.
The negative powers come only from differentiating
\(y=e^{Q/t^2}\). For dependence on an open domain, replace this
exponential above the lip by
\(t^M\chi((Q-Q_{\rm lip})/t^2)\), where
\(Q_{\rm lip}=Mt^2\log t\), \(\chi(z)=e^z\) for
\(z\le1\), and \(\chi\) is smooth, positive and constant for
\(z\ge2\). This changes no first-exit orbit, preserves
\eqref{neg:core-mixed-estimate}, and keeps the denominator positive.

The two comparison exit cuts from \eqref{neg:action-equation}
have opposite action differences. Starting at the positive lip,
\(Q\) decreases, stays a fixed amount below the lip near
\(\xi=0\), and increases strictly on the negative side away from
zero. The comparison cuts consequently give a unique actual first
exit. At that exit,
\(Q_\xi=-r/P_\Lambda(-r)+o(1)\) is bounded away from zero.
The differentiated integral equation for \eqref{neg:thin-core},
followed by this implicit exit equation, therefore has uniform
derivatives through total order five. With the two lip maps it gives
\begin{equation}\label{neg:compact-mixed-convergence}
\begin{split}
 &\|S_{t,\lambda,k}-S_{0,\lambda,k}\|_{C_b^3}
  +\|\partial_\lambda S_{t,\lambda,k}
                -\partial_\lambda S_{0,\lambda,k}\|_{C_b^3}\\
 &\hspace{20mm}\le C\{t^2(1+|\log t|)+t^{M-10}\},
 \qquad S_{0,\lambda,k}(b)=-r^2/2 .
\end{split}
\end{equation}
No derivative with respect to \(t\) at zero is asserted.

The exact energy and even \(P_{0,k}\) give
\(S_{t,0,k}=S_{0,0,k}=\mathrm{id}\). At \(\lambda=k=0\),
\eqref{neg:action-equation} gives
\[
 \partial_{\Lambda_1}S_{0,0,0}=-2x^3/3=U_0(b),\qquad
 \partial_{\Lambda_2}S_{0,0,0}=2x^5/5=V_0(b).
\]
Its derivatives are smooth on the compact root-free block. The
identity on the coefficient ideal permits the exact division
\[
 S_{t,\lambda,k}-\mathrm{id}
      =\int_0^1\partial_\lambda S_{t,v\lambda,k}\,\lambda\,dv .
\]
Using \eqref{neg:compact-mixed-convergence} in this integral, and
then the smoothness of the limiting action, gives
\begin{equation}\label{neg:relative-core}
 \|S_{t,\lambda,k}-\mathrm{id}-\Lambda_1U_0-\Lambda_2V_0\|_{C^3}
  \le C\rho_{\rm id}
          \{t^2(1+|\log t|)+\rho_{\rm id}+|k|\} .
\end{equation}
Thus division by \(\rho_{\rm id}\) precedes the limit; the even
coefficient \(k\) creates no additional row.

The same construction bounds \(S\) in \(C^4\). Apply
Lemma~\ref{phys:lem-c3-core-transport} with
\eqref{neg:single-tail-estimate} to obtain
\begin{equation}\label{neg:composition-estimate}
 \|\Pi_{t,\Lambda}-S_{t,\lambda,k}\|_{C^3}
                       \le Ct^2\rho_{\rm id} .
\end{equation}
Only the core needs the extra phase derivative for evaluation at
\(\phi_+^{-1}\). Combining with \eqref{neg:relative-core} proves
\eqref{neg:relative-estimate}.
\end{proof}

The outgoing tail must first be estimated in its own intrinsic label. If
an outgoing graph is written as \(\eta(\xi,S(b))\), then
\[
 \partial_b^2\eta(\xi,S(b))
    =\eta_{BB}(\xi,S(b))(S')^2+\eta_B(\xi,S(b))S'' .
\]
The second term need not decay on that tail. Estimating the map in
\(B\) first and then composing retains this contribution.

\subsection{Zeros of the completed displacement}

The action transition has a curvature formula useful beyond the
small coefficient neighborhood. Differentiating
\eqref{neg:action-equation}, with \(x=\sqrt{-2b}\), gives
\begin{equation}\label{neg:action-curvature}
 S_\Lambda'=
     \frac{P_\Lambda(-r)}{P_\Lambda(x)}>0,
 \qquad
 \frac{S_\Lambda''}{S_\Lambda'}
   =\frac{x+r}{P_\Lambda(x)}
                 \left(3\Lambda_2-\frac{\Lambda_1}{xr}\right),
 \quad S_\Lambda=-r^2/2 .
\end{equation}
The terms containing \(k\) cancel. Moreover
\begin{equation}\label{neg:entry-exit-monotonicity}
 \frac{d(xr)}{db}
      =-\frac r x-\frac{xP_\Lambda(-r)}{rP_\Lambda(x)}<0 .
\end{equation}
Thus, away from \(\lambda=0\), every possible curvature zero is
simple and has the same crossing orientation.

\begin{corollary}\label{neg:cor-rootfree}
On a compact family satisfying the section, tail, and root-free
margins above, assume \(\Lambda_2\ge0\) and that the physical
coefficients satisfy \(k\to0\) whenever
\(\rho_{\rm id}\to0\), uniformly on the family. For example,
the condition \(|k|^3\le C\Lambda_2^2\) suffices.
There is a common \(t_0>0\) such that every actual completed
displacement \eqref{neg:physical-displacement}, for
\(0<t\le t_0\), has at most three isolated zeros on its input
interval, counted with multiplicity.

The same conclusion holds for connected coefficient-dependent
intervals \(J_\Lambda\) if their coefficient--label graph has
compact closure and the enlarged input and output intervals, both
exit margins, and all actual transition neighborhoods have the
same uniform margins.
\end{corollary}
\begin{proof}
At \(\rho_{\rm id}=0\), the exact first integral and the separate
normalized tail limits give \(\Pi=\mathrm{id}\). The displacement
is therefore affine: it has at most one zero unless it vanishes
identically, in which case it has no isolated zero. This coefficient
boundary is understood as a normalized limit; actual raw passages
have \(\Lambda_2>0\).

Near the coefficient ideal, put \(x=\sqrt{-2b}\). The limiting
rows satisfy
\[
 U_0''=-2/x,\qquad V_0''=6x,\qquad
 \frac{(\theta_1U_0+\theta_2V_0)''}{2/x}
                   =-\theta_1+3\theta_2x^2 .
\]
If the last expression vanishes, compactness of the positive
\(x\)-interval and \(\theta_1^2+\theta_2^2=1\) give a
positive lower bound on \(\theta_2\). Its derivative in \(b\)
is \(-6\theta_2\). When \(\theta_2\) is smaller than a fixed
positive constant, the expression is instead bounded away from
zero. This proves a uniform alternative of no curvature zero or
one simple downward crossing. The relative \(C^3\) estimate
preserves that alternative after a common shrink of
\(t,\rho_{\rm id},k\).

On the complementary compact coefficient annulus, the tail estimate
still applies as \(p=t^2\Lambda_1\), \(s=t^2\Lambda_2\),
\(\gamma=t^2k\) tend uniformly to zero. The strict core margin permits
the same lip and first-exit argument, giving absolute \(C^3\)
convergence to \(S_\Lambda\) without division through the ideal. Hence
\(\|[\Pi-S_\Lambda]\|_{C^3/\operatorname{Aff}}\to0\).
Equations \eqref{neg:action-curvature} and
\eqref{neg:entry-exit-monotonicity} give the same uniform crossing
alternative on this compact annulus. Finitely overlapping the two
coefficient regions gives one bound on \(t\).

Lemma~\ref{sharp:upper-curvature-count} gives at most one curvature
zero and three displacement zeros since \(D''=-mc_-\Pi''\).
The affine relation \eqref{neg:physical-displacement} preserves
multiplicity. Compactness of the coefficient--label graph supplies
the same constants fibrewise on \(J_\Lambda\).
\end{proof}
